% Research proposition; to be read conditional on Lemmas \ref{lem:reduction},
% \ref{lem:valuations}, \ref{lem:model} and the proof of Theorem \ref{thm:lower}.
\subsection*{Calibration of clipped first-passage mass}
\label{subsec:clipped-calibration}
Throughout this subsection let $c\ge-1$ be odd with $3\nmid c$, fix
$\beta=1/200$, and set $\gamma_0=2/\log(4/3)$.
For $X>1$ put
\[
 W_X=(X^{1-\beta},X]\cap\N,\qquad X_1=X^{1-2\beta},\qquad
 H_X=(\log X)^{\varepsilon_0/2}X.
\]
For an odd positive integer $m$, let $\Phi_X(m)$ be the sum, counted
with multiplicity, of the elements of $W_X$ in Syracuse blocks
$1,\ldots,\tau_{X_1}(m)$, and put $\Phi_X(m)=0$ if
$\tau_{X_1}(m)=\infty$. For $w\in\N$, let $w^\flat$ be its odd part.
\begin{proposition}[Sharp calibration after clipping]
\label{prop:clipped-calibration}
Assume Lemmas~\ref{lem:reduction}, \ref{lem:valuations}, and
\ref{lem:model} with the uniformities used in the proof of
Theorem~\ref{thm:lower}. Then the logarithmic mean
\[
 C_X:=\lim_{N\to\infty}\frac1{\log N}
 \sum_{1\le w<N}\frac{\min\{\Phi_X(w^\flat),H_X\}}{w}
\]
exists for every sufficiently large $X$, and
\begin{equation}\label{eq:clipped-calibration}
 \boxed{C_X=\gamma_0 X+o(X)\qquad(X\to\infty).}
\end{equation}
\end{proposition}
\begin{proof}
For fixed $X$, the first-passage invariance checked in Step~2 of
the proof of Theorem~\ref{thm:lower} holds for
$\varphi_X(m)=\min\{\Phi_X(m),H_X\}/H_X$.
Lemma~\ref{lem:reduction} therefore gives the existence of $C_X$
and, with $x_0=X^{1+\beta}$, $u=x_0^{\lambda_0}$,
$v=u^\alpha$, and $\mathbf N\sim\mathbf{Log}_{u,v}$,
\begin{equation}\label{eq:calib-window}
 C_X=\mathbb E\min\{\Phi_X(\mathbf N),H_X\}
  +O\bigl(H_X(\log u)^{-\varepsilon_0}\bigr)
 =\mathbb E\min\{\Phi_X(\mathbf N),H_X\}+o(X).
\end{equation}
The proof of Theorem~\ref{thm:lower}, Step~3, with the residue
condition removed, shows
\begin{equation}\label{eq:calib-lower}
 \mathbb E\min\{\Phi_X(\mathbf N),H_X\}\ge\gamma_0X-o(X).
\end{equation}
Here one must apply the argument to the \emph{total} block sum;
summing separately clipped residue contributions would not suffice.

We prove the reverse inequality. Couple $\mathbf N$ to
$U\sim\mathrm{Unif}[s,s+L]$ and independent geometric variables
$g_j$ as in Step~3, with $s=\log u$, $L=\log(v/u)$, and
$\varepsilon=X^{-1/2}$. Let $G$ be the good coupling event there;
for some $\eta>0$ it satisfies $\mathbb P(G^c)=O(X^{-\eta})$.
For $j\ge1$ and $0\le i<g_j$, write
$z_{j,i}=U-D_j+i\log2$ as in Lemma~\ref{lem:model}.
The comparison in Step~3 gives uniformly for
$j\le\tau_{X_1}(\mathbf N)$ on $G$,
\[
 \left|\log(2^i\operatorname{Syr}^j(\mathbf N))-z_{j,i}\right|
 \le C\varepsilon,
\]
for all $0\le i<g_j$; in particular the valuations agree up to
this stopping time. (The relative-error formula
\eqref{eq:ideal} uses only preceding ideal iterates, which stay
above $X_1/3$ even for the final block.) Consequently, with
$\ell_-=(1-\beta)\log X-C\varepsilon$,
$\ell_+=\log X+C\varepsilon$, and
\[
 Z^+=\sum_{j\ge1}\sum_{i=0}^{g_j-1}e^{z_{j,i}}
       \mathbf 1[\ell_-<z_{j,i}\le\ell_+],
\]
we have the pointwise upper comparison
\begin{equation}\label{eq:calib-upper-coupling}
 \Phi_X(\mathbf N)\le (1+C\varepsilon)Z^+\quad\text{on }G.
\end{equation}

Let $V(t)$ be the expected number of unmarked model points
$y_{j,i}=D_j-i\log2$ below $t$, over all $j\ge1$ and
$0\le i<g_j$. Summing the three residue classes in the proof of
Lemma~\ref{lem:model}, equation~\eqref{eq:V} for $n=1$, gives
$V(t)=\gamma_0t+O(1)$ for $t\ge0$. Thus the same exact
integration identity as in the proof of Lemma~\ref{lem:model}(a)
gives
\[
 \begin{aligned}
 \mathbb EZ^+
 &=\frac1L\int_{\ell_-}^{\ell_+}e^z
             \bigl(V(s+L-z)-V(s-z)\bigr)\,dz\\
 &=\gamma_0\bigl(e^{\ell_+}-e^{\ell_-}\bigr)
       +O\left(\frac{e^{\ell_+}}L\right)
  =\gamma_0 X+o(X),
 \end{aligned}
\]
because $s-\ell_+>0$, $L\asymp\log X$,
$e^{\ell_+}=X(1+O(\varepsilon))$, and
$e^{\ell_-}=O(X^{1-\beta})$.
Since $\min\{\Phi_X,H_X\}\le H_X$ always,
\[
 \mathbb E\min\{\Phi_X(\mathbf N),H_X\}
 \le(1+C\varepsilon)\mathbb EZ^+
     +H_X\mathbb P(G^c)
 =\gamma_0X+o(X).
\]
Together with \eqref{eq:calib-lower} and
\eqref{eq:calib-window}, this proves the claim.
\end{proof}

\begin{corollary}[Two nonnegative obstructions to the mean upper bound]
\label{cor:two-tails}
Write $S_X=\sum_{a\le X}a h(a)$, and define
\[
\begin{aligned}
 F_X&=\lim_{N\to\infty}\frac1{\log N}
   \sum_{1\le w<N}\frac{\Phi_X(w^\flat)}w,\\
 R_X&=\lim_{N\to\infty}\frac1{\log N}
   \sum_{1\le w<N}\frac1w
   \left(\sum_{\substack{a\le X\\a\in\operatorname{orb}(w)}}a
   -\Phi_X(w^\flat)\right),
\end{aligned}
\]
where the limits exist by Lemma~\ref{lem:reduction},
Proposition~\ref{prop:potential}(a), and finite linearity.
Put
\[
 T_X=F_X-C_X
 =\lim_{N\to\infty}\frac1{\log N}
     \sum_{1\le w<N}\frac{(\Phi_X(w^\flat)-H_X)_+}{w}\ge0.
\]
Then
\[
 \boxed{S_X=C_X+T_X+R_X,\qquad T_X\ge0,\quad R_X\ge0.}
\]
In particular, conditional on the stated lemmas,
\[
 \limsup_{X\to\infty}\frac{S_X}{X}\le\gamma_0
 \quad\Longleftrightarrow\quad T_X=o(X)\ \text{and}\ R_X=o(X).
\]
\end{corollary}
\begin{proof}
For fixed $X$, the first-passage blocks do not repeat points,
by Step~1 of Theorem~\ref{thm:lower}. Thus for every $w$,
$0\le\Phi_X(w^\flat)\le
\sum_{a\le X,\ a\in\operatorname{orb}(w)}a$.
Lemma~\ref{lem:reduction} applies to the bounded function
$\Phi_X/B_X$, where $B_X=\sum_{a\in W_X}a$, since the first-passage
invariance of Step~2 holds for all sufficiently large passage
thresholds. Hence $F_X$ exists. Proposition~\ref{prop:potential}(a)
and finite linearity show that the logarithmic mean of the full
orbit-point sum is $S_X$. Subtracting proves that $R_X$ exists,
$R_X\ge0$, and $S_X=F_X+R_X$.
Finally $F_X-C_X$ is the mean of $(\Phi_X-H_X)_+$, so $T_X\ge0$. The claimed equivalence follows
from Proposition~\ref{prop:clipped-calibration} and the fact that
both tails are nonnegative.
\end{proof}

\paragraph{A concrete next target: uniform integrability.}
A sufficient (currently unproved) condition for the clipping tail to
vanish is the existence of $\theta>0$ and $K<\infty$ such that,
uniformly in $X$,
\[
 \limsup_{N\to\infty}\frac1{\log N}
 \sum_{1\le w<N}\frac1w
 \left(\frac{\Phi_X(w^\flat)}X\right)^{1+\theta}
 \le K.
\]
Indeed, $(t-H_X)_+\le t^{1+\theta}/H_X^\theta$
for $t\ge0$, whence
\[
 0\le T_X\le K X (\log X)^{-\varepsilon_0\theta/2}=o(X).
\]
This condition says nothing about the distinct residual term $R_X$.
